\setLTRbibitems

\begin{thebibliography}{999}

\addcontentsline{toc}{chapter}{منابع و مآخذ}

\small

\singlespacing

\bibitem{ref1} World Health Organization. (2021). Suicide worldwide in 2019: Global health estimates. Geneva: WHO.

\bibitem{ref2} Centers for Disease Control and Prevention (CDC). (2022). Leading Causes of Death Reports, 2020. National Vital Statistics System.

\bibitem{ref3} Turecki, G., \& Brent, D. A. (2016). Suicide and suicidal behaviour. The Lancet, 387(10024), 1227-1239.

\bibitem{ref4} Blease, C., Torous, J., \& Hagmann, C. (2020). Ethical challenges in the use of digital data for mental health research. The Lancet Digital Health, 2(8), e386-e387.

\bibitem{ref5} Glenn, C. R., et al. (2020). Adolescent suicide risk and the developmental context. Annual Review of Clinical Psychology, 16, 397-422.

\bibitem{ref6} Niederkrotenthaler, T., et al. (2020). Association of suicidal ideation and behavior with social media use. JAMA Network Open, 3(10), e2021398.

\bibitem{ref7} Twenge, J. M., et al. (2018). Increases in depression and suicide among US adolescents. Clinical Psychological Science, 6(1), 3-17.

\bibitem{ref8} Coppersmith, G., et al. (2018). Natural language processing in mental health. ACM Computing Surveys, 51(5), 1-36.

\bibitem{ref9} Gunnarsson, K. F., et al. (2021). Linguistic markers of suicidal ideation: A systematic review. Journal of Affective Disorders, 295, 1205-1217.

\bibitem{ref10} Shatte, A. B. R., et al. (2019). Machine learning in mental health: A scoping review. JMIR Mental Health, 6(11), e13316.

\bibitem{ref11} Ji, S., Pan, S., Li, X., Cambria, E., Long, G., \& Huang, Z. (2021). Suicidal ideation detection: A review of machine learning methods and applications. IEEE Transactions on Computational Social Systems, 8(1), 214–226.

\bibitem{ref12} Heckler, W. F., de Carvalho, J. V., \& Barbosa, J. L. V. (2021). Machine learning for suicidal ideation identification: A systematic literature review. Computers in Human Behavior, 120, 106807.

\bibitem{ref13} Leveraging Reddit for suicidal ideation detection: A review of machine learning and natural language processing techniques. (2022).

\bibitem{ref14} Sheng, Z. (2024). Suicidal ideation detection on social media using machine learning: A review. Applied and Computational Engineering, 71, 58–63.

\bibitem{ref15} McMahan, B., Moore, E., Ramage, D., Hampson, S., \& Arcas, B. A. y. (2017). Communication-efficient learning of deep networks from decentralized data. Proceedings of the 20th International Conference on Artificial Intelligence and Statistics, 54, 1273–1282.

\bibitem{ref16} Rieke, N., Hancox, J., Li, W., et al. (2020). The future of digital health with federated learning. NPJ Digital Medicine, 3, 119.

\bibitem{ref17} Kairouz, P., McMahan, H. B., Avent, B., et al. (2021). Advances and open problems in federated learning. Foundations and Trends in Machine Learning, 14(1–2), 1–210.

\bibitem{ref18} Khalil, S. S., Abdelsalam, N., Spruit, M., et al. (2024). Federated learning for privacy- preserving depression detection with multilingual language models in social media posts. Frontiers in Digital Health.

\bibitem{ref19} Paillier, P. (1999). Public-key cryptosystems based on composite degree residuosity classes. In Advances in Cryptology—EUROCRYPT ’99, 223–238.

\bibitem{ref20} Gentry, C. (2009). A fully homomorphic encryption scheme. Stanford University.

\bibitem{ref21} Wood, A., Najarian, K., \& Kahrobaei, D. (2021). Homomorphic encryption for machine learning in medicine and bioinformatics. ACM Computing Surveys, 53(4).

\bibitem{ref22} A systematic review of homomorphic encryption and its contributions in healthcare industry. (2022).

\bibitem{ref23} Lee, C. H., Lim, K. H., \& Eswaran, S. (2025). A comprehensive survey on secure healthcare data processing with homomorphic encryption: Attacks and defenses. Discover Public Health, 22, 137.

\bibitem{ref24} Khurana, D., Koli, A., Khatter, K., \& Singh, S. (2021). Natural Language Processing: State of The Art, Current Trends and Challenges.

\bibitem{ref25} Chomsky, N. (1965). Aspects of the Theory of Syntax. Cambridge, MA: MIT Press.

\bibitem{ref26} Rospocher, M., van Erp, M., Vossen, P., Fokkens, A., Aldabe, I., Rigau, G., Soroa, A., Ploeger, T., \& Bogaard, T. (2016). Building event-centric knowledge graphs from news. Web Semantics: Science, Services and Agents on the World Wide Web, In Press.

\bibitem{ref27} Shemtov, H. (1997). Ambiguity management in natural language generation. Stanford University.

\bibitem{ref28} Hutchins, W. J. (1986). Machine translation: past, present, future (p. 66). Chichester: Ellis Horwood.

\bibitem{ref29} Hutchins, W. J. (Ed.). (2000). Early years in machine translation: memoirs and biographies of pioneers (Vol. 97). John Benjamins Publishing.

\bibitem{ref30} Green Jr., B. F., Wolf, A. K., Chomsky, C., \& Laughery, K. (1961, May). Baseball: an automatic question-answerer. In Papers presented at the May 9-11, 1961, western joint IRE-AIEE-ACM computer conference (pp. 219-224). ACM.

\bibitem{ref31} Woods, W. A. (1978). Semantics and quantification in natural language question answering. Advances in computers, 17, 1-87.

\bibitem{ref32} Hendrix, G. G., Sacerdoti, E. D., Sagalowicz, D., \& Slocum, J. (1978). Developing a natural language interface to complex data. ACM Transactions on Database Systems (TODS), 3(2), 105-147.

\bibitem{ref33} Wilks, Y., \& Fass, D. (1992). "The SRI Core Language Engine." In Natural Language Processing Systems. Springer-Verlag.

\bibitem{ref34} Kamp, H., \& Reyle, U. (1993). From Discourse to Logic: Introduction to Model- theoretic Semantics of Natural Language, Formal Logic and Discourse Representation Theory. Kluwer Academic Publishers.

\bibitem{ref35} Briscoe, T., Grover, C., Boguraev, B., \& Carroll, J. (1987). "The Alvey Natural Language Tools." In Proceedings of the 3rd Conference of the European Chapter of the Association for Computational Linguistics, pp. 139-148.

\bibitem{ref36} Young, S., \& Chase, L. (1998). "Speech Recognition and Understanding: The DARPA Program." In Proceedings of the IEEE, 86(5), pp. 1267-1283.

\bibitem{ref37} Sundheim, B., \& Chinchor, N. (1993). "Survey of the Message Understanding Conferences." In Proceedings of the Workshop on Human Language Technology, pp. 56-60.

\bibitem{ref38} Kobsa, A., \& Wahlster, W. (Eds.). (1989). User Models in Dialog Systems. Springer- Verlag.

\bibitem{ref39} Cohen, P. R., Morgan, J., \& Pollack, M. E. (Eds.). (1990). Intentions in Communication. MIT Press.

\bibitem{ref40} McKeown, K. R. (1985). Text Generation: Using Discourse Strategies and Focus Constraints to Generate Natural Language Text. Cambridge University Press.

\bibitem{ref41} Small, S. L., Cottrell, G. W., \& Tanenhaus, M. K. (Eds.). (1988). Lexical Ambiguity Resolution: Perspectives from Psycholinguistics, Neuropsychology, and Artificial Intelligence. Morgan Kaufmann.

\bibitem{ref42} Manning, C. D., \& Schütze, H. (1999). Foundations of Statistical Natural Language Processing. MIT Press.

\bibitem{ref43} Mani, I., \& Maybury, M. T. (Eds.). (1999). Advances in Automatic Text Summarization. MIT Press.

\bibitem{ref44} Kilgarriff, A., \& Grefenstette, G. (2003). Introduction to the special issue on the web as corpus. Computational Linguistics, 29(3), 333-347.

\bibitem{ref45} Haddi, E., Liu, X., \& Shi, Y. (2013). The role of text pre-processing in sentiment analysis. Procedia Computer Science, 17, 26-32.

\bibitem{ref46} Kaur, A., \& Gupta, V. (2020). A comprehensive survey on text preprocessing techniques for natural language processing. ACM Computing Surveys, 52(4), 1-36.

\bibitem{ref47} Aljuaid, H., et al. (2020). Importance of preprocessing in Arabic text classification. IEEE Access, 8, 205230-205243.

\bibitem{ref48} Loper, E., \& Bird, S. (2002). NLTK: The natural language toolkit. Proceedings of the ACL-02 Workshop on Effective Tools and Methodologies for Teaching NLP, 63-70.

\bibitem{ref49} Kowsari, K., et al. (2019). Text classification algorithms: A survey. Information, 10(4),

\bibitem{ref50} Eisner, B., et al. (2019). Emoji-based sentiment analysis. Proceedings of the 2019 Conference on Empirical Methods in Natural Language Processing, 2345-2350.

\bibitem{ref51} Fridkin, S., \& M. (2020). Challenges in multilingual NLP for low-resource languages. Journal of Artificial Intelligence Research, 68, 245-298.

\bibitem{ref52} Shamsfard, M., \& S. (2019). Persian text normalization: A comprehensive approach. Language Resources and Evaluation, 53(2), 267-295.

\bibitem{ref53} Karimi, S., et al. (2020). Challenges in Persian text processing. ACM Transactions on Asian and Low-Resource Language Information Processing, 19(3), 1-25.

\bibitem{ref54} Wilbur, J., \& S. (2019). Stop word removal in natural language processing. Journal of Computational Linguistics, 45(1), 127-152.

\bibitem{ref55} Vijayarani, S., \& J. (2020). Performance analysis of text preprocessing techniques. International Journal of Advanced Computer Science and Applications, 11(6), 345-352.

\bibitem{ref56} Manning, C. D., \& Schütze, H. (1999). Foundations of statistical natural language processing. MIT Press.

\bibitem{ref57} Jurafsky, D., \& Martin, J. H. (2023). Speech and language processing (3rd ed.). Stanford University.

\bibitem{ref58} Coppersmith, G., et al. (2018). Natural language processing in mental health. ACM Computing Surveys, 51(5), 1-36.

\bibitem{ref59} Loper, E., \& Bird, S. (2002). NLTK: The natural language toolkit. Proceedings of the ACL-02 Workshop on Effective Tools and Methodologies for Teaching NLP, 63-70.

\bibitem{ref60} Bird, S., Klein, E., \& Loper, E. (2009). Natural language processing with Python. O'Reilly Media.

\bibitem{ref61} Kaur, A., \& Gupta, V. (2020). A comprehensive survey on text preprocessing techniques for natural language processing. ACM Computing Surveys, 52(4), 1-36.

\bibitem{ref62} Sennrich, R., Haddow, B., \& Birch, A. (2016). Neural machine translation of rare words with subword units. Proceedings of the 54th Annual Meeting of the Association for Computational Linguistics, 1715-1725.

\bibitem{ref63} Liu, Y., et al. (2021). Efficient text preprocessing for large-scale NLP. Journal of Big Data, 8(1), 1-18.

\bibitem{ref64} Fridkin, S., \& M. (2020). Challenges in multilingual NLP for low-resource languages. Journal of Artificial Intelligence Research, 68, 245-298.

\bibitem{ref65} Gunnarsson, K. F., et al. (2021). Linguistic markers of suicidal ideation: A systematic review. Journal of Affective Disorders, 295, 1205-1217.

\bibitem{ref66} Vijayarani, S., \& J. (2020). Performance analysis of text preprocessing techniques. International Journal of Advanced Computer Science and Applications, 11(6), 345-352.

\bibitem{ref67} Zhang, C., et al. (2021). Privacy-preserving machine learning: A survey. ACM Computing Surveys, 54(7), 1-36.

\bibitem{ref68} Haddi, E., Liu, X., \& Shi, Y. (2013). The role of text pre-processing in sentiment analysis. Procedia Computer Science, 17, 26-32.

\bibitem{ref69} Shamsfard, M., \& S. (2019). Persian text normalization: A comprehensive approach. Language Resources and Evaluation, 53(2), 267-295.

\bibitem{ref70} Karimi, S., et al. (2020). Challenges in Persian text processing. ACM Transactions on Asian and Low-Resource Language Information Processing, 19(3), 1-25.

\bibitem{ref71} Aljuaid, H., et al. (2020). Importance of preprocessing in Arabic text classification. IEEE Access, 8, 205230-205243.

\bibitem{ref72} Eskandari, M., \& M. (2021). Spelling correction and spacing in Persian language. Journal of Language Resources and Evaluation, 55(1), 45-68.

\bibitem{ref73} Salton, G., \& Buckley, C. (1988). Term-weighting approaches in automatic text retrieval. Information Processing \& Management, 24(5), 513-523.

\bibitem{ref74} Kowsari, K., et al. (2019). Text classification algorithms: A survey. Information, 10(4),

\bibitem{ref75} Devlin, J., et al. (2019). BERT: Pre-training of deep bidirectional transformers for language understanding. Proceedings of NAACL-HLT, 4171-4186.

\bibitem{ref76} Manning, C. D., \& Schütze, H. (1999). Foundations of statistical natural language processing. MIT Press.

\bibitem{ref77} Bengio, Y., et al. (2003). A neural probabilistic language model. Journal of Machine Learning Research, 3, 1137-1155.

\bibitem{ref78} Mikolov, T., et al. (2013). Efficient estimation of word representations in vector space. arXiv:1301.3781.

\bibitem{ref79} Coppersmith, G., et al. (2018). Natural language processing in mental health. ACM Computing Surveys, 51(5), 1-36.

\bibitem{ref80} Reimers, N., \& Gurevych, I. (2019). Sentence-BERT: Sentence embeddings using Siamese BERT-networks. Proceedings of EMNLP-IJCNLP, 3982-3992.

\bibitem{ref81} Vaswani, A., et al. (2017). Attention is all you need. Advances in Neural Information Processing Systems, 30, 5998-6008.

\bibitem{ref82} Devlin, J., et al. (2019). BERT: Pre-training of deep bidirectional transformers for language understanding. Proceedings of NAACL-HLT, 4171-4186.

\bibitem{ref83} Wang, W., et al. (2020). MiniLM: Deep self-attention distillation for task-agnostic compression of pre-trained transformers. Advances in Neural Information Processing Systems, 33, 5776-5788.

\bibitem{ref84} Loper, E., \& Bird, S. (2002). NLTK: The natural language toolkit. Proceedings of the ACL-02 Workshop, 63-70.

\bibitem{ref85} Pedregosa, F., et al. (2011). Scikit-learn: Machine learning in Python. Journal of Machine Learning Research, 12, 2825-2830.

\bibitem{ref86} Thakur, N., et al. (2021). Augmented SBERT: Data augmentation method for improving bi-encoders. Proceedings of EACL, 141-152.

\bibitem{ref87} Salton, G., \& Buckley, C. (1988). Term-weighting approaches in automatic text retrieval. Information Processing \& Management, 24(5), 513-523.

\bibitem{ref88} Jones, K. S. (1972). A statistical interpretation of term specificity and its application in retrieval. Journal of Documentation, 28(1), 11-21.

\bibitem{ref89} Ramos, J. (2003). Using TF-IDF to determine word relevance in document queries. Proceedings of the First Instructional Conference on Machine Learning, 242, 29-48.

\bibitem{ref90} Liu, Y., et al. (2019). RoBERTa: A robustly optimized BERT pretraining approach. arXiv:1907.11692.

\bibitem{ref91} Peters, M. E., et al. (2018). Deep contextualized word representations. Proceedings of NAACL-HLT, 2227-2237.

\bibitem{ref92} He, Q., et al. (2021). Transformer-based models for suicide risk detection. Journal of Biomedical Informatics, 118, 103793.

\bibitem{ref93} Strubell, E., et al. (2019). Energy and policy considerations for deep learning in NLP. Proceedings of ACL, 3645-3650.

\bibitem{ref94} Conneau, A., et al. (2020). Unsupervised cross-lingual representation learning at scale. Proceedings of ACL, 8440-8451.

\bibitem{ref95} Manning, C. D., et al. (2008). Introduction to information retrieval. Cambridge University Press.

\bibitem{ref96} Vijayarani, S., \& J. (2020). Performance analysis of text preprocessing techniques. International Journal of Advanced Computer Science and Applications, 11(6), 345-352.

\bibitem{ref97} Aljuaid, H., et al. (2020). Importance of preprocessing in Arabic text classification. IEEE Access, 8, 205230-205243.

\bibitem{ref98} Kowsari, K., et al. (2019). Text classification algorithms: A survey. Information, 10(4),

\bibitem{ref99} Haddi, E., Liu, X., \& Shi, Y. (2013). The role of text pre-processing in sentiment analysis. Procedia Computer Science, 17, 26-32.

\bibitem{ref100} Kaur, A., \& Gupta, V. (2020). A comprehensive survey on text preprocessing techniques for natural language processing. ACM Computing Surveys, 52(4), 1-36.

\bibitem{ref101} Zhang, C., et al. (2021). Privacy-preserving machine learning: A survey. ACM Computing Surveys, 54(7), 1-36.

\bibitem{ref102} Reimers, N., \& Gurevych, I. (2020). Making monolingual sentence embeddings multilingual using knowledge distillation. Proceedings of EMNLP, 4512-4525.

\bibitem{ref103} Breiman, L. (2001). Random Forests. Machine Learning, 45(1), 5-32.

\bibitem{ref104} Ho, T. K. (1995). Random decision forests. Proceedings of the 3rd International Conference on Document Analysis and Recognition, 278-282.

\bibitem{ref105} Amit, Y., \& Geman, D. (1997). Shape quantization and recognition with randomized trees. Neural Computation, 9(7), 1545-1588.

\bibitem{ref106} Biau, G., \& Scornet, E. (2016). A random forest guided tour. Test, 25(2), 197-227.

\bibitem{ref107} Fawagreh, K., Gaber, M. M., \& Elyan, E. (2014). Random forests: from early developments to recent advancements. Systems Science \& Control Engineering, 2(1), 602-609.

\bibitem{ref108} Hastie, T., Tibshirani, R., \& Friedman, J. (2009). The Elements of Statistical Learning: Data Mining, Inference, and Prediction. Springer.

\bibitem{ref109} Dietterich, T. G. (2000). Ensemble methods in machine learning. Multiple Classifier Systems, 1-15.

\bibitem{ref110} Provost, F., \& Fawcett, T. (2013). Data Science for Business: What You Need to Know about Data Mining and Data-Analytic Thinking. O'Reilly Media.

\bibitem{ref111} He, H., \& Garcia, E. A. (2009). Learning from imbalanced data. IEEE Transactions on Knowledge and Data Engineering, 21(9), 1263-1284.

\bibitem{ref112} Zhang, H. (2004). The optimality of Naive Bayes. Proceedings of FLAIRS 2004.

\bibitem{ref113} GHOST: Adjusting the Decision Threshold to Handle Imbalanced Data in Machine Learning. (2021). PubMed (NIH).

\bibitem{ref114} Risk-Aware Calibration and Interpretable Baselines for Mental Health Text Classification. (2026). IEEE Xplore.

\bibitem{ref115} Ding et al. (2024). Hyper-parameters setting for machine learning models. PeerJ Computer Science, 10, e2248.

\bibitem{ref116} Chawla, N. V., Bowyer, K. W., Hall, L. O., \& Kegelmeyer, W. P. (2002). SMOTE: Synthetic minority over-sampling technique. Journal of Artificial Intelligence Research, 16, 321-357.

\bibitem{ref117} Fawcett, T. (2006). An introduction to ROC analysis. Pattern Recognition Letters, 27(8), 861-874.

\bibitem{ref118} Davis, J., \& Goadrich, M. (2006). The relationship between Precision-Recall and ROC curves. Proceedings of the 23rd International Conference on Machine Learning, 233-240.

\bibitem{ref119} Sokolova, M., Japkowicz, N., \& Szpakowicz, S. (2006). Beyond accuracy, F-score and ROC: a family of discriminant measures for performance evaluation. AI 2006: Advances in Artificial Intelligence, 1015-1021.

\bibitem{ref120} Machine Learning-Based Classification of Suicidal Ideation in Tweets. (2025). IEEE Xplore.

\bibitem{ref121} Bazrafshan, M., \& Sayehmiri, K. (2024). Predicting suicidal behavior outcomes: an analysis of key factors and machine learning models. BMC Psychiatry, 24, 123.

\bibitem{ref122} Cawley, G. C., \& Talbot, N. L. C. (2010). On over-fitting in model selection and subsequent selection bias in performance evaluation. Journal of Machine Learning Research, 11, 2079-2107.

\bibitem{ref123} López, V., Fernández, A., García, S., Palade, V., \& Herrera, F. (2013). An insight into classification with imbalanced data: Empirical results and current trends on using data intrinsic characteristics. Information Sciences, 250, 113-141.

\bibitem{ref124} Safavian, S. R., \& Landgrebe, D. (1991). A survey of decision tree classifier methodology. IEEE Transactions on Systems, Man, and Cybernetics, 21(3), 660-674.

\bibitem{ref125} Hunt, E. B., Marin, J., \& Stone, P. J. (1966). Experiments in Induction. Academic Press.

\bibitem{ref126} Quinlan, J. R. (1986). Induction of decision trees. Machine Learning, 1(1), 81-106.

\bibitem{ref127} Quinlan, J. R. (1993). C4.5: Programs for Machine Learning. Morgan Kaufmann.

\bibitem{ref128} Breiman, L., Friedman, J., Stone, C. J., \& Olshen, R. A. (1984). Classification and Regression Trees. CRC Press.

\bibitem{ref129} Pedregosa, F., et al. (2011). Scikit-learn: Machine learning in Python. Journal of Machine Learning Research, 12, 2825-2830.

\bibitem{ref130} Mitchell, T. M. (1997). Machine Learning. McGraw-Hill.

\bibitem{ref131} Han, J., Kamber, M., \& Pei, J. (2011). Data Mining: Concepts and Techniques. Morgan Kaufmann.

\bibitem{ref132} Rokach, L., \& Maimon, O. (2005). Decision trees. In Data Mining and Knowledge Discovery Handbook, 165-192.

\bibitem{ref133} Witten, I. H., Frank, E., \& Hall, M. A. (2011). Data Mining: Practical Machine Learning Tools and Techniques. Morgan Kaufmann.

\bibitem{ref134} Loh, W. Y. (2011). Classification and regression trees. Wiley Interdisciplinary Reviews: Data Mining and Knowledge Discovery, 1(1), 14-23.

\bibitem{ref135} Caruana, R., \& Niculescu-Mizil, A. (2006). An empirical comparison of supervised learning algorithms. Proceedings of the 23rd International Conference on Machine Learning, 161-168.

\bibitem{ref136} Wu, X., et al. (2008). Top 10 algorithms in data mining. Knowledge and Information Systems, 14(1), 1\lr{--}37.

\bibitem{ref137} Dietterich, T. G. (1995). Overfitting and undercomputing in machine learning. ACM Computing Surveys, 27(3), 326-327.

\bibitem{ref138} Esposito, F., Malerba, D., \& Semeraro, G. (1997). A comparative analysis of methods for pruning decision trees. IEEE Transactions on Pattern Analysis and Machine Intelligence, 19(5), 476-491.

\bibitem{ref139} Zadrozny, B., \& Elkan, C. (2001). Obtaining calibrated probability estimates from decision trees and naive Bayesian classifiers. Proceedings of the 18th International Conference on Machine Learning, 609-616.

\bibitem{ref140} Provost, F., \& Fawcett, T. (2013). Data Science for Business: What You Need to Know about Data Mining and Data-Analytic Thinking. O'Reilly Media.

\bibitem{ref141} He, H., \& Garcia, E. A. (2009). Learning from imbalanced data. IEEE Transactions on Knowledge and Data Engineering, 21(9), 1263-1284.

\bibitem{ref142} Zhang, H. (2004). The optimality of Naive Bayes. Proceedings of FLAIRS 2004.

\bibitem{ref143} Liu, W., Chawla, S., Cieslak, D. A., \& Chawla, N. V. (2011). A robust decision tree algorithm for imbalanced data sets. Proceedings of the 2011 SIAM International Conference on Data Mining, 766-777.

\bibitem{ref144} Risk-Aware Calibration and Interpretable Baselines for Mental Health Text Classification. (2026). IEEE Xplore.

\bibitem{ref145} Sokolova, M., Japkowicz, N., \& Szpakowicz, S. (2006). Beyond accuracy, F-score and ROC: a family of discriminant measures for performance evaluation. AI 2006: Advances in Artificial Intelligence, 1015-1021.

\bibitem{ref146} Fawcett, T. (2006). An introduction to ROC analysis. Pattern Recognition Letters, 27(8), 861-874.

\bibitem{ref147} GHOST: Adjusting the Decision Threshold to Handle Imbalanced Data in Machine Learning. (2021). PubMed (NIH).

\bibitem{ref148} López, V., Fernández, A., García, S., Palade, V., \& Herrera, F. (2013). An insight into classification with imbalanced data: Empirical results and current trends on using data intrinsic characteristics. Information Sciences, 250, 113-141.

\bibitem{ref149} Bazrafshan, M., \& Sayehmiri, K. (2024). Predicting suicidal behavior outcomes: an analysis of key factors and machine learning models. BMC Psychiatry, 24, 123.

\bibitem{ref150} Haghish, E. F., et al. (2023). Reconsidering False Positives in Machine Learning Binary Classification Models of Suicidal Behavior. Current Psychology, 43, 10117-10121.

\bibitem{ref151} Alasmari, M. J. (2022). Ensemble Learning and XGBoost: A Comprehensive Review. SACM - United Kingdom.

\bibitem{ref152} Gulati, J. (2025). Algorithm Showdown: Logistic Regression vs. Random Forest vs. XGBoost on Imbalanced Data. MachineLearningMastery.com.

\bibitem{ref153} Text-Based Mental Health Detection Using Stacking Ensemble Approach with Hybrid Machine Learning and Deep Learning Models. (2026). IEEE Xplore.

\bibitem{ref154} Artificial intelligence approach for predicting suicide-related behaviour in emergency departments. (2026). Scientific Reports (Nature).

\bibitem{ref155} Mienye, I. D., \& Sun, Y. (2022). A Survey of Ensemble Learning: Concepts, Algorithms, Applications, and Prospects. IEEE Access, 10, 99129-99149.

\bibitem{ref156} ParaBoostForest Hybrid Parallel Boosting for Imbalanced Learning. (2025). GitHub.

\bibitem{ref157} Wang, H. (2026). UEDU: Mental health distress signal detection in student writing. GitHub.

\bibitem{ref158} Anchor rules from XGBoost model for suicide prediction. (2024). NIH/PMC.

\bibitem{ref159} Chen, T., \& Guestrin, C. (2016). XGBoost: A Scalable Tree Boosting System. Proceedings of the 22nd ACM SIGKDD.

\bibitem{ref160} XGBoost training and threshold optimization. (2026). Imperial College London GitLab.

\bibitem{ref161} Early detection of mental health on social media using a hybrid Bi-LSTM–XGBoost model. (2026). Scientific Reports (Springer), 16, 20582.

\bibitem{ref162} Garuda – Garba Rujukan Digital. (2026). Systematic Literature Review of Naive Bayes Algorithm for Data Classification in Indonesia.

\bibitem{ref163} Garuda – Garba Rujukan Digital. (2025). Twitter Sentiment Analysis on the Iran-Israel Conflict Using the Naïve Bayes Classification Algorithm.

\bibitem{ref164} National Institutes of Health (NIH). (2018). Potential use of text classification tools as signatures of suicidal behavior: A proof-of-concept study using Virginia Woolf’s personal writings. PLoS ONE, 13(10), e0204820.

\bibitem{ref165} IEEE Xplore. (2026). Comparative Analysis of Naïve Bayes Algorithms for Lecturer Performance Classification.

\bibitem{ref166} MDPI. (2024). Detecting Suicidal Ideations in Online Forums with Textual and Psycholinguistic Features. Applied Sciences, 14(21), 9911.

\bibitem{ref167} Journal of Computer Science and Information Technology. (2025). Twitter Sentiment Analysis on the Iran-Israel Conflict Using the Naïve Bayes Classification Algorithm.

\bibitem{ref168} OUCI. (2023). Ensemble Learning for Email Classification using Voting Classifier. Expert Systems with Applications.

\bibitem{ref169} Journal of Nnamdi Azikiwe University. (2026). An Improved Multinomial Naïve Bayes Algorithm with Laplace Smoothing for News Classification.

\bibitem{ref170} GitHub. (2023). Suicide Detection using Naive Bayes Classification.

\bibitem{ref171} GitHub. (2025). Ensemble Learning Voting Classifier using Naive Bayes.

\bibitem{ref172} Virascience. (2015). A Survey Paper On Naive Bayes Classifier For Multi-Feature Based Text Mining.

\bibitem{ref173} Nauki i Tekhniki. (2024). Review of methods for semantic text classification.

\bibitem{ref174} Czech Academy of Sciences. (2023). Spam Detection Using Ensemble Learning.

\bibitem{ref175} Stack Overflow. (2015). Naive Bayes Text Classification Laplace Smoothing.

\bibitem{ref176} Zhang, H. (2004). The optimality of Naive Bayes. Proceedings of the Seventeenth International Florida Artificial Intelligence Research Society Conference (FLAIRS 2004).

\bibitem{ref177} Eskandari, M., \& M. (2021). Spelling correction and spacing in Persian language. Journal of Language Resources and Evaluation, 55(1), 45-68.

\bibitem{ref178} ElGamal, T. (1985). A Public Key Cryptosystem and a Signature Scheme Based on Discrete Logarithms. IEEE Transactions on Information Theory, 31(4), 469–472.

\bibitem{ref179} Paillier, P. (1999). Public-Key Cryptosystems Based on Composite Degree Residuosity Classes. EUROCRYPT 1999, LNCS 1592, 223–238.

\bibitem{ref180} Benaloh, J. (1994). Dense Probabilistic Encryption. Selected Areas in Cryptography

\bibitem{ref181} Damgård, I., \& Jurik, M. (2001). A Generalisation, a Simplification and Some Applications of Paillier’s Probabilistic Public-Key System. PKC 2001, LNCS 1992, 119\lr{--}136.

\bibitem{ref182} Rivest, R. L., Shamir, A., \& Adleman, L. (1978). A Method for Obtaining Digital Signatures and Public-Key Cryptosystems. Communications of the ACM, 21(2), 120\lr{--}126.

\bibitem{notebook-models} Ranjbar Yazdi, S. Corrected model-comparison notebook, \texttt{finding-the-best-model-fixed.ipynb}. Source and measured outputs archived with this thesis. \url{https://github.com/Sajjad-s-presentations/My-master-s-thesis/blob/presentation/Notebooks/finding-the-best-model-fixed.ipynb}.

\bibitem{notebook-federated} Ranjbar Yazdi, S. Federated XGBoost experiments with 10, 20, 50 and 100 clients. Archived executed notebook and extracted result tables. \url{https://github.com/Sajjad-s-presentations/My-master-s-thesis/blob/presentation/Notebooks/federated_xgboost_10_20_50_100.ipynb}.

\bibitem{shamir} Shamir, A. (1979). How to Share a Secret. Communications of the ACM, 22(11), 612\lr{--}613. doi:10.1145/359168.359176.

\bibitem{nist-gcm} Dworkin, M. (2007). Recommendation for Block Cipher Modes of Operation: Galois/Counter Mode (GCM) and GMAC. NIST Special Publication 800-38D. doi:10.6028/NIST.SP.800-38D.

\bibitem{ed25519} Josefsson, S., \& Liusvaara, I. (2017). Edwards-Curve Digital Signature Algorithm (EdDSA). RFC 8032. doi:10.17487/RFC8032.

\bibitem{bonawitz} Bonawitz, K., et al. (2017). Practical Secure Aggregation for Privacy-Preserving Machine Learning. Proceedings of the 2017 ACM SIGSAC Conference on Computer and Communications Security, 1175\lr{--}1191. doi:10.1145/3133956.3133982.

\end{thebibliography}
